\label{sec:discussion}
\head{What changes for people running coordination alarms.} The baseline a monitor is validated against decides what its alarms mean. Validated against independent agents, today's monitors report fleets as coordinating: at the Hugging Face swarm's scale, hundreds of thousands of alerts with no communication behind them. The monitors that avoid this do not detect transmission. The fix is cheap and needs no model: compare agents with agents from \emph{other runs of the same deployment}, which inherits whatever the runs share (model, harness, docs, queue, outages) without anyone having to list it; evaluation fleets already consist of such runs. This gives OpenAI's ``cross-run coordination'' trigger [2] an operational definition with a known false-alarm rate, and population-level control [16, 17] the cost number it lacks.

\head{Limitations.} M4 needs replicate runs whose common causes recur (\texttt{fleet\_drift} shows it fails otherwise) and detects shared provenance, not direct edges, and only channels that leave a trace in actions. Ground truth and M4's signal share one trace, so a simple rarity rule matches M4 here. One task, one shock, one model family and an instructed swarm limit generality; at fleet scale the pairwise level needs a correction over pairs. Full list and dual use: Appendix~\ref{app:limits}.

\head{Future work.} Evaluation fleets with replicate runs and several model families; a population-level error rate; gateway read logs, which turn provenance into edges; channels that leave no action trace [10, 19].

\head{Dated predictions.} P1 (70\%): by 30 Jun 2027 a public incident analysis will report agent ``coordination'' that a later analysis attributes mainly to common cause. P2 (60\%): by 31 Dec 2027 no frontier lab will have published its coordination trigger's false-alarm rate on a no-channel control. P3 (60\%): the next multi-agent collusion benchmark after 8 Nov 2026 will lack a same-model, shared-shock, no-channel control.
